\section{Discussion}
\label{sec:discussion}

\subsection{Choosing an engine}
\label{sec:guide}

\Cref{tab:decision} condenses the results into guidance. Its central message is that
the right engine depends on three properties of the deployment: how heterogeneous the request
lengths are, which latency the SLO is bound by, and how much KV-cache headroom the model leaves.
\todo{finalise the table once \cref{sec:results} is written}

\begin{table}[t]
  \caption{Engine choice by deployment situation.}
  \label{tab:decision}
  \small
  \begin{tabular}{@{}p{3.3cm}p{3.0cm}l@{}}
    \toprule
    Situation & Prefer & Basis \\
    \midrule
    Uniform requests, steady load & \trt & \rqref{1} \\
    Mixed lengths, TTFT-bound SLO & \vllm & \rqref{1} \\
    Smooth streaming required & \sglang / \trt & \rqref{1} \\
    Model barely fits in KV memory & \vllm, more memory & \rqref{1}, \rqref{4} \\
    Slot-bound deployment & raise \texttt{max\_num\_seqs} & \rqref{1} \\
    Long prompts & chunked prefill on \trt only & \rqref{2} \\
    Multi-turn chat & prefix caching on any engine & \rqref{2} \\
    N-gram speculation & \vllm only & \rqref{3} \\
    \bottomrule
  \end{tabular}
\end{table}

\subsection{Lessons for benchmarking LLM serving}
\label{sec:lessons}

Several measurement choices that look innocuous changed our conclusions. We report them because
they apply to any LLM serving benchmark.

\paragraph{Measure TTFT from arrival, not from dispatch.}
A client that caps in-flight requests and starts the TTFT clock at dispatch silently excludes
the time requests spend queued at the client---the classic coordinated-omission
problem~\cite{tene2015latency}. In our data this hidden wait reached a 95th percentile of over
200\,s for the most overloaded engine, so dispatch-based TTFT understated the gap between
engines precisely where it was largest. Open-loop arrival-based measurement~\cite{schroeder2006open}
avoids this; for closed-loop multi-turn workloads, dispatch-based TTFT is the correct choice.

\paragraph{Count timeouts; do not drop them.}
Discarding requests that time out removes an overloaded engine's worst latencies and can invert
rankings. Censoring them at the timeout value and counting them as SLO misses keeps the
comparison honest.

\paragraph{Finite traces make latency non-monotonic in load.}
With a fixed number of requests, a low arrival rate stretches the trace over minutes, so late
arrivals meet a KV cache already full of long-running generations, while at a high rate all
requests arrive before the cache fills. In our data the 95th-percentile TTFT of one
configuration was 185\,s at 2 requests/s and 0.3\,s at 16 requests/s; the cost moved from the
first token to mid-stream stalls. Reporting latency against position in the run, or using
duration-based traces, exposes this effect.

\paragraph{Identically named engine metrics are not comparable.}
Engines report prefix-cache hit rates over different windows (cumulative, last batch, or once
per run), and the gap between streamed chunks shrinks or grows with the number of tokens per
chunk under speculative decoding. We therefore compare engines only on client-side metrics and
use engine-reported values to explain them. TPOT, unlike TBT, is robust to multi-token chunks.

\paragraph{Prompt formatting affects load, not just quality.}
Sending unformatted prompts to instruction-tuned models changed how many tokens they generated
and, for some models, produced replies with zero tokens. Both change the load an engine
experiences; a benchmark should either apply each model's chat template or report output-length
distributions alongside latency.

\paragraph{Per-token speed is a poor proxy for service quality.}
Under the SLOs we study, goodput was decided almost entirely by TTFT: requests that missed only
the TPOT target were at most 2\% in every configuration. Benchmarks that report tokens per second
at a comfortable load measure the property that least often limits real deployments.
